\documentclass[12pt,a4paper]{article}
\usepackage[UTF8]{ctex}
\usepackage{geometry}
\usepackage{graphicx}
\usepackage{fancyhdr}
\usepackage{listings}
\usepackage{xcolor}
\usepackage{float}
\usepackage{longtable}
\usepackage{booktabs}
\usepackage{hyperref}
\usepackage{enumitem}

\geometry{left=2.5cm,right=2.5cm,top=2.5cm,bottom=2.5cm}

% 代码样式设置
\lstset{
    basicstyle=\small\ttfamily,
    keywordstyle=\color{blue},
    commentstyle=\color{green!50!black},
    stringstyle=\color{red},
    numberstyle=\tiny\color{gray},
    numbers=left,
    numbersep=5pt,
    frame=single,
    breaklines=true,
    breakatwhitespace=false,
    tabsize=4,
    showspaces=false,
    showstringspaces=false,
    captionpos=b,
    xleftmargin=2em,
    xrightmargin=2em,
}

\lstdefinestyle{cpp}{
    language=C++,
    morekeywords={std,uint32_t,SOCKET,sockaddr_in,atomic,mutex,thread}
}

\hypersetup{
    colorlinks=true,
    linkcolor=black,
    filecolor=blue,
    urlcolor=blue,
    citecolor=cyan,
}

\pagestyle{fancy}
\fancyhf{}
\fancyhead[C]{浙江大学实验报告}
\fancyfoot[C]{\thepage}

\begin{document}

% 封面
\begin{titlepage}
    \centering
    \vspace*{1cm}

    \includegraphics[width=0.3\textwidth]{ZJU.png}

    \vspace{1.5cm}

    {\Huge \textbf{本科实验报告}}

    \vspace{2cm}

    \begin{tabular}{rl}
        \textbf{课程名称：} & 计算机网络基础 \\[0.5cm]
        \textbf{实验名称：} & 基于Socket接口实现自定义协议通信 \\[0.5cm]
        \textbf{姓\qquad 名：} & 姜快快 \\[0.5cm]
        \textbf{学\qquad 院：} & 计算机学院 \\[0.5cm]
        \textbf{专\qquad 业：} & 计算机科学与技术 \\[0.5cm]
        \textbf{学\qquad 号：} & 3210104669 \\[0.5cm]
        \textbf{指导教师：} & \qquad\qquad\qquad \\[0.5cm]
    \end{tabular}

    \vfill

    {\large \today}

\end{titlepage}

\newpage
\thispagestyle{fancy}

% 实验信息表
\begin{center}
    {\Large \textbf{浙江大学实验报告}}
\end{center}

\vspace{0.5cm}

\noindent
\begin{tabular}{ll}
    \textbf{实验名称：} & 基于Socket接口实现自定义协议通信 \\[0.3cm]
    \textbf{实验类型：} & 编程实验 \\[0.3cm]
    \textbf{同组学生：} & \qquad\qquad\qquad\qquad\qquad \\[0.3cm]
    \textbf{实验地点：} & 计算机网络实验室 \\[0.3cm]
\end{tabular}

\section{实验目的}
掌握Socket编程接口编写基本的网络应用软件。

\section{实验内容}

根据自定义的协议规范，使用Socket编程接口编写基本的网络应用软件。

\begin{itemize}[leftmargin=2em]
    \item 掌握C语言形式的Socket编程接口用法，能够正确发送和接收网络数据包
    \item 开发一个客户端，实现人机交互界面和与服务器的通信
    \item 开发一个服务端，实现并发处理多个客户端的请求
    \item 程序界面不做要求，使用命令行或最简单的窗体即可
\end{itemize}

\subsection{功能要求}

\subsubsection{运输层协议}
运输层协议采用TCP。

\subsubsection{客户端功能}
客户端采用交互菜单形式，用户可以选择以下功能：

\begin{enumerate}
    \item \textbf{连接：}请求连接到指定地址和端口的服务端
    \item \textbf{断开连接：}断开与服务端的连接
    \item \textbf{获取时间：}请求服务端给出当前时间
    \item \textbf{获取名字：}请求服务端给出其机器的名称
    \item \textbf{活动连接列表：}请求服务端给出当前连接的所有客户端信息（编号、IP地址、端口等）
    \item \textbf{发消息：}请求服务端把消息转发给对应编号的客户端，该客户端收到后显示在屏幕上
    \item \textbf{退出：}断开连接并退出客户端程序
\end{enumerate}

\subsubsection{服务端功能}
服务端接收到客户端请求后，根据客户端传过来的指令完成特定任务：

\begin{itemize}
    \item 向客户端传送服务端所在机器的当前时间
    \item 向客户端传送服务端所在机器的名称
    \item 向客户端传送当前连接的所有客户端信息
    \item 将某客户端发送过来的内容转发给指定编号的其他客户端
\end{itemize}

\subsubsection{编程要求}
\begin{itemize}
    \item 采用异步多线程编程模式，正确处理多个客户端同时连接、同时发送消息的情况
    \item 本实验涉及到网络数据包发送部分不能使用任何的Socket封装类，只能使用最底层的C语言形式的Socket API
    \item 服务端监听端口使用学号后四位：\textbf{4669}
\end{itemize}

\section{主要仪器设备}

\begin{itemize}
    \item 联网的PC机
    \item Windows操作系统
    \item MinGW C++编译器（支持C++17标准）
    \item Winsock2库
\end{itemize}

\section{操作方法与实验步骤}

\subsection{协议设计}

本实验采用自定义应用层协议，报文格式如下：

\begin{center}
\begin{tabular}{|c|c|c|}
\hline
\textbf{字段} & \textbf{长度} & \textbf{说明} \\
\hline
type & 4字节 & 消息类型（网络字节序） \\
\hline
request\_id & 4字节 & 请求ID（网络字节序） \\
\hline
body\_length & 4字节 & 消息体长度（网络字节序） \\
\hline
body & 变长 & 消息体内容（UTF-8编码） \\
\hline
\end{tabular}
\end{center}

报头共12字节，所有整数采用网络字节序（大端序）。消息类型定义如下：

\begin{center}
\begin{tabular}{|l|c|}
\hline
\textbf{消息类型} & \textbf{类型值} \\
\hline
请求时间 & 1 \\
请求名字 & 2 \\
请求客户端列表 & 3 \\
请求发送消息 & 4 \\
\hline
响应时间 & 101 \\
响应名字 & 102 \\
响应客户端列表 & 103 \\
转发消息 & 104 \\
消息转发确认 & 105 \\
问候消息 & 106 \\
错误响应 & 199 \\
\hline
\end{tabular}
\end{center}

\subsection{客户端编写步骤}

客户端采用多线程模式，包括主线程和接收线程。

\subsubsection{主线程职责}
\begin{enumerate}
    \item 运行初始化，调用 \texttt{WSAStartup()} 和 \texttt{socket()}
    \item 显示菜单功能，列出7个选项
    \item 等待用户选择并执行相应操作
    \item 处理接收线程放入消息队列的数据
    \item 显示响应消息和转发消息
\end{enumerate}

\subsubsection{接收线程职责}
\begin{enumerate}
    \item 循环调用 \texttt{recv()}，接收服务端数据
    \item 拼接并解析完整数据包
    \item 将响应消息或指示消息放入线程安全消息队列
    \item 收到主线程退出通知后结束
\end{enumerate}

\subsection{服务端编写步骤}

服务端采用多线程模式，主线程负责接受连接，每个客户端由独立的工作线程处理。

\subsubsection{主线程职责}
\begin{enumerate}
    \item 初始化Winsock，创建监听Socket
    \item 绑定学号后四位端口（4669）
    \item 调用 \texttt{listen()} 设置连接等待队列
    \item 循环调用 \texttt{accept()} 接受客户端连接
    \item 为每个客户端创建处理线程
\end{enumerate}

\subsubsection{工作线程职责}
\begin{enumerate}
    \item 发送问候消息（可选）
    \item 循环调用 \texttt{recv()}，接收客户端请求
    \item 根据请求类型执行相应操作：
    \begin{itemize}
        \item 获取时间：调用 \texttt{time()} 并返回
        \item 获取名字：调用 \texttt{GetComputerNameW()} 并返回
        \item 获取客户端列表：返回所有活动客户端信息
        \item 转发消息：将消息转发给目标客户端
    \end{itemize}
    \item 连接断开或异常时清理资源
\end{enumerate}

\section{实验数据记录和处理}

\subsection{客户端和服务器端框架图}

% 预留位置插入框架图
\begin{figure}[H]
    \centering
    % TODO: 插入框架图
    \fbox{\parbox{0.8\textwidth}{\centering\vspace{5cm}[预留位置：客户端和服务端框架流程图]\vspace{5cm}}}
    \caption{客户端和服务端框架流程图}
\end{figure}

\subsection{客户端初始运行菜单}

% 预留位置插入菜单截图
\begin{figure}[H]
    \centering
    % TODO: 插入截图
    \fbox{\parbox{0.8\textwidth}{\centering\vspace{3cm}[预留位置：客户端初始菜单截图]\vspace{3cm}}}
    \caption{客户端初始运行菜单}
\end{figure}

客户端启动后显示7个菜单选项，包括连接、断开连接、获取时间、获取名字、列出活动客户端、发送消息和退出。未连接时只能选择连接和退出功能。

\subsection{客户端接收数据子线程关键代码}

客户端接收线程的核心代码如下（详见附录A）：

\begin{lstlisting}[style=cpp,caption={客户端接收线程循环}]
void receiverLoop() {
    Packet packet;
    while (!stopReceiver.load() && recvPacket(clientSocket, packet)) {
        enqueue(std::move(packet));
        packet = Packet{};
    }
    if (!stopReceiver.load()) {
        disconnectedEvent = true;
        enqueue({0, 0, "server connection closed"});
    }
}
\end{lstlisting}

接收线程持续调用 \texttt{recvPacket()} 函数，该函数先读取12字节报头，再根据报头中的长度字段读取消息体，确保接收完整数据包。

\subsection{服务器初始运行界面}

% 预留位置插入服务器截图
\begin{figure}[H]
    \centering
    % TODO: 插入截图
    \fbox{\parbox{0.8\textwidth}{\centering\vspace{3cm}[预留位置：服务器初始运行界面截图]\vspace{3cm}}}
    \caption{服务器初始运行界面}
\end{figure}

服务器启动后显示监听端口4669，并提示按Ctrl+C停止服务器。

\subsection{服务器客户端处理子线程关键代码}

服务器工作线程的核心代码如下（详见附录B）：

\begin{lstlisting}[style=cpp,caption={服务器工作线程循环}]
void clientWorker(const std::shared_ptr<Client>& client) {
    sendToClient(client, {kGreeting, 0, "Connected to Lab2 server"});

    Packet request;
    while (client->connected.load() && recvPacket(client->socket, request)) {
        switch (request.type) {
        case kRequestTime:
            // 处理时间请求
            break;
        case kRequestName:
            // 处理名字请求
            break;
        case kRequestList:
            // 处理列表请求
            break;
        case kRequestMessage:
            // 处理消息转发
            break;
        }
    }
    removeClient(client);
}
\end{lstlisting}

每个客户端由独立线程处理，使用该客户端的Socket进行收发，不使用监听Socket。

\subsection{客户端选择连接功能}

% 预留位置插入连接截图
\begin{figure}[H]
    \centering
    % TODO: 插入截图
    \fbox{\parbox{0.8\textwidth}{\centering\vspace{4cm}[预留位置：客户端连接功能截图（客户端和服务端）]\vspace{4cm}}}
    \caption{客户端连接功能}
\end{figure}

客户端输入服务器IP地址（如127.0.0.1），使用固定端口4669连接。连接成功后服务端显示客户端编号和地址信息。

\subsection{客户端选择获取时间功能}

% 预留位置插入时间截图
\begin{figure}[H]
    \centering
    % TODO: 插入截图
    \fbox{\parbox{0.8\textwidth}{\centering\vspace{4cm}[预留位置：获取时间功能截图（客户端和服务端）]\vspace{4cm}}}
    \caption{获取服务器时间}
\end{figure}

\textbf{服务端处理代码：}

\begin{lstlisting}[style=cpp,caption={服务端处理时间请求}]
case kRequestTime: {
    log("[server] client #" + std::to_string(client->id) + " requested server time");
    const std::time_t now = std::time(nullptr);
    std::tm localTime{};
    localtime_s(&localTime, &now);
    std::ostringstream timeText;
    timeText << std::put_time(&localTime, "%Y-%m-%d %H:%M:%S");
    sendToClient(client, {kResponseTime, request.requestId, timeText.str()});
    break;
}
\end{lstlisting}

\subsection{客户端选择获取名字功能}

% 预留位置插入名字截图
\begin{figure}[H]
    \centering
    % TODO: 插入截图
    \fbox{\parbox{0.8\textwidth}{\centering\vspace{4cm}[预留位置：获取名字功能截图（客户端和服务端）]\vspace{4cm}}}
    \caption{获取服务器名字}
\end{figure}

\textbf{服务端处理代码：}

\begin{lstlisting}[style=cpp,caption={服务端处理名字请求}]
case kRequestName: {
    log("[server] client #" + std::to_string(client->id) + " requested server name");
    wchar_t name[MAX_COMPUTERNAME_LENGTH + 1]{};
    DWORD length = MAX_COMPUTERNAME_LENGTH + 1;
    if (!GetComputerNameW(name, &length)) {
        sendToClient(client, {kResponseError, request.requestId, "GetComputerName failed"});
        break;
    }
    const int size = WideCharToMultiByte(CP_UTF8, 0, name, static_cast<int>(length), nullptr, 0, nullptr, nullptr);
    std::string value(size, '\0');
    WideCharToMultiByte(CP_UTF8, 0, name, static_cast<int>(length), value.data(), size, nullptr, nullptr);
    sendToClient(client, {kResponseName, request.requestId, value});
    break;
}
\end{lstlisting}

\subsection{客户端选择获取客户端列表功能}

% 预留位置插入列表截图
\begin{figure}[H]
    \centering
    % TODO: 插入截图
    \fbox{\parbox{0.8\textwidth}{\centering\vspace{4cm}[预留位置：获取客户端列表功能截图（客户端和服务端）]\vspace{4cm}}}
    \caption{获取活动客户端列表}
\end{figure}

\textbf{服务端处理代码：}

\begin{lstlisting}[style=cpp,caption={服务端生成客户端列表}]
std::string makeClientList() {
    std::vector<std::shared_ptr<Client>> snapshot;
    {
        std::lock_guard<std::mutex> lock(clientsMutex);
        for (const auto& [id, client] : clients) {
            if (client->connected.load()) snapshot.push_back(client);
        }
    }
    std::sort(snapshot.begin(), snapshot.end(),
              [](const auto& a, const auto& b) { return a->id < b->id; });
    std::ostringstream result;
    result << snapshot.size() << '\n';
    for (const auto& client : snapshot) {
        char ip[INET_ADDRSTRLEN]{};
        inet_ntop(AF_INET, &client->address.sin_addr, ip, sizeof(ip));
        result << client->id << ' ' << ip << ' ' << ntohs(client->address.sin_port) << '\n';
    }
    return result.str();
}
\end{lstlisting}

\subsection{客户端选择发送消息功能}

% 预留位置插入发送消息截图
\begin{figure}[H]
    \centering
    % TODO: 插入截图
    \fbox{\parbox{0.8\textwidth}{\centering\vspace{3cm}[预留位置：发送消息功能截图（发送客户端）]\vspace{3cm}}}
    \caption{发送消息的客户端}
\end{figure}

\begin{figure}[H]
    \centering
    % TODO: 插入截图
    \fbox{\parbox{0.8\textwidth}{\centering\vspace{3cm}[预留位置：接收消息功能截图（接收客户端）]\vspace{3cm}}}
    \caption{接收消息的客户端}
\end{figure}

\begin{figure}[H]
    \centering
    % TODO: 插入截图
    \fbox{\parbox{0.8\textwidth}{\centering\vspace{3cm}[预留位置：服务端转发消息截图（可选）]\vspace{3cm}}}
    \caption{服务端转发消息}
\end{figure}

\textbf{服务端消息转发代码：}

\begin{lstlisting}[style=cpp,caption={服务端转发消息}]
case kRequestMessage: {
    const auto separator = request.body.find('\n');
    int targetId = 0;
    if (!parsePositiveInt(request.body.substr(0, separator), targetId, 2147483647)) {
        sendToClient(client, {kResponseError, request.requestId, "invalid target client number"});
        break;
    }
    const std::string message = request.body.substr(separator + 1);

    std::shared_ptr<Client> target;
    {
        std::lock_guard<std::mutex> lock(clientsMutex);
        const auto it = clients.find(targetId);
        if (it != clients.end()) target = it->second;
    }

    if (!target || !target->connected.load()) {
        sendToClient(client, {kResponseError, request.requestId, "target client is not connected"});
        break;
    }

    std::ostringstream forwarded;
    forwarded << "from #" << client->id << ' ' << addressText(client->address) << "\n" << message;
    if (sendToClient(target, {kForwardMessage, 0, forwarded.str()})) {
        sendToClient(client, {kResponseMessageAck, request.requestId, "message forwarded"});
        log("[server] forwarded #" + std::to_string(client->id) + " -> #" + std::to_string(targetId));
    }
    break;
}
\end{lstlisting}

\textbf{客户端发送消息代码：}

\begin{lstlisting}[style=cpp,caption={客户端发送消息}]
case 6: {
    if (!clientListFetched.load()) {
        std::cout << "Retrieve the active client list with option 5 before sending a message.\n";
        break;
    }
    const std::string target = readLineWithEvents("Target client number (use option 5 first): ");
    int targetId = 0;
    try { targetId = std::stoi(target); } catch (...) { targetId = 0; }
    const std::string message = readLineWithEvents("Message: ");
    waitForResponse(sendRequest(kRequestMessage, std::to_string(targetId) + "\n" + message));
    break;
}
\end{lstlisting}

\textbf{客户端接收消息代码：}

\begin{lstlisting}[style=cpp,caption={客户端处理转发消息}]
case kForwardMessage:
    std::cout << "[forwarded message]\n" << packet.body << "\n";
    break;
\end{lstlisting}

\section{实验结果与分析}

\subsection{客户端是否需要调用bind操作？}

客户端不需要调用 \texttt{bind()} 操作。客户端的源端口由操作系统自动分配。在调用 \texttt{connect()} 时，如果Socket未绑定端口，操作系统会从临时端口范围（通常是49152-65535）中选择一个可用端口。

每次调用 \texttt{connect()} 时，如果使用新创建的Socket，客户端的端口通常会变化，因为操作系统会分配一个新的临时端口。如果复用同一个Socket，端口保持不变。

\subsection{服务端listen和accept之间的断点测试}

假设在服务端调用 \texttt{listen()} 和调用 \texttt{accept()} 之间设置断点，此时客户端调用 \texttt{connect()} 可以立即成功。这是因为 \texttt{listen()} 函数已经创建了连接等待队列，客户端的连接请求会被放入该队列中，即使服务端尚未调用 \texttt{accept()}。

\texttt{listen()} 的第二个参数指定了等待队列的最大长度（本程序使用 \texttt{SOMAXCONN}）。只要队列未满，客户端连接就能成功，但实际的数据传输要等到服务端调用 \texttt{accept()} 后才能进行。

\subsection{服务器如何区分不同客户端的数据包？}

服务器在同一个监听端口接收多个客户端的数据，通过不同的Socket句柄来区分数据包属于哪个客户端。

具体过程如下：
\begin{enumerate}
    \item 服务器主线程在监听Socket上调用 \texttt{accept()}，每接受一个新连接，\texttt{accept()} 返回一个新的Socket句柄
    \item 每个客户端对应一个唯一的Socket句柄和一个独立的工作线程
    \item 工作线程只在自己的Socket句柄上调用 \texttt{recv()}，因此接收到的数据必然来自对应的客户端
    \item 服务器维护一个客户端表，记录每个客户端的编号、Socket句柄、IP地址和端口
\end{enumerate}

本程序中，每个客户端由 \texttt{Client} 结构体表示，包含唯一ID和Socket句柄，存储在全局的 \texttt{clients} 映射表中。

\subsection{客户端主动断开后的TCP连接状态}

客户端主动断开连接后，TCP连接进入 \texttt{TIME\_WAIT} 状态。该状态通常保持 2MSL（Maximum Segment Lifetime）时间，在Windows系统中通常是120秒（2分钟）。

使用 \texttt{netstat -an} 命令可以观察到：
\begin{itemize}
    \item 主动关闭方（客户端）进入 \texttt{TIME\_WAIT} 状态
    \item 被动关闭方（服务端）先进入 \texttt{CLOSE\_WAIT}，然后进入 \texttt{LAST\_ACK}，最后关闭
\end{itemize}

\texttt{TIME\_WAIT} 状态的目的是：
\begin{enumerate}
    \item 确保最后的ACK能够到达对方
    \item 确保旧连接的数据包在网络中消失，避免与新连接混淆
\end{enumerate}

\subsection{客户端异常退出时服务器的TCP状态变化}

客户端断网后异常退出，服务器的TCP连接不会立即检测到断开，连接可能长时间保持在 \texttt{ESTABLISHED} 状态。这是因为TCP是双向独立的，一方断开不会立即通知另一方。

服务器检测连接是否有效的方法：
\begin{enumerate}
    \item \textbf{调用recv()检测：}当连接断开时，\texttt{recv()} 返回0或SOCKET\_ERROR，这是本程序采用的方法
    \item \textbf{设置SO\_KEEPALIVE选项：}启用TCP保活机制，定期发送探测包
    \item \textbf{应用层心跳：}定期发送心跳消息，超时未收到响应则认为连接断开
    \item \textbf{设置接收超时：}使用SO\_RCVTIMEO选项，超时后主动检测连接状态
\end{enumerate}

本程序中，工作线程在 \texttt{recvPacket()} 中调用 \texttt{recv()}，如果客户端异常断开，\texttt{recv()} 会返回错误，工作线程检测到后调用 \texttt{removeClient()} 清理资源。

\section{讨论与心得}

\subsection{实验过程中遇到的困难}

\begin{enumerate}
    \item \textbf{TCP粘包和拆包问题：}TCP是流式协议，一次 \texttt{send()} 可能被分成多个TCP段，一次 \texttt{recv()} 可能接收到多个消息或部分消息。通过实现 \texttt{sendAll()} 和 \texttt{recvExact()} 函数解决，确保发送和接收指定字节数。

    \item \textbf{多线程并发控制：}多个线程同时访问客户端列表和Socket句柄需要加锁保护。使用 \texttt{std::mutex} 保护客户端列表的增删改查，使用每个客户端独立的 \texttt{sendMutex} 防止同时发送时数据交错。

    \item \textbf{字节序转换：}网络传输使用大端序（网络字节序），而Intel x86架构使用小端序（主机字节序）。所有整数字段需要使用 \texttt{htonl()} 和 \texttt{ntohl()} 进行转换。

    \item \textbf{客户端主线程阻塞问题：}如果主线程等待用户输入时使用阻塞式 \texttt{std::cin}，无法及时处理接收线程的消息。通过实现 \texttt{readLineWithEvents()} 函数，使用 \texttt{\_kbhit()} 和 \texttt{\_getch()} 非阻塞读取用户输入，同时定期调用 \texttt{drainMessages()} 处理接收到的消息。

    \item \textbf{资源清理和线程安全退出：}服务端退出时需要正确关闭监听Socket、通知所有工作线程退出、等待线程结束、关闭所有客户端Socket。通过 \texttt{std::atomic} 标志和 \texttt{shutdown()} 函数实现优雅退出。
\end{enumerate}

\subsection{得到的经验教训}

\begin{enumerate}
    \item Socket编程必须处理部分发送和部分接收，不能假定一次系统调用对应一个完整消息。
    \item 多线程程序需要仔细设计并发控制，避免数据竞争和死锁。
    \item 网络程序需要处理各种异常情况：连接断开、超时、对方异常退出等。
    \item 使用自定义协议时，报文格式设计要考虑扩展性和健壮性，包括版本号、长度字段、类型字段等。
    \item 测试时应该使用多个客户端同时连接，模拟并发场景，才能发现潜在的并发问题。
\end{enumerate}

\subsection{对本实验的建议}

\begin{enumerate}
    \item 建议提供更详细的协议设计示例，减少学生在协议设计上花费的时间。
    \item 建议提供跨平台的实验环境，或者同时提供Windows和Linux的示例代码。
    \item 建议增加性能测试要求，如测试最大并发连接数、消息吞吐量等。
    \item 建议增加异常处理测试，如突然断网、服务端崩溃等场景。
\end{enumerate}

\newpage

\appendix

\section{附录A：客户端完整源代码}

\lstinputlisting[style=cpp,caption={user.cpp}]{repo/Lab2/user.cpp}

\newpage

\section{附录B：服务端完整源代码}

\lstinputlisting[style=cpp,caption={server.cpp}]{repo/Lab2/server.cpp}

\end{document}
